\chapter{System architecture}
\label{ch:architecture}

Yap has two tiers: a native mobile client and a hosted Supabase project. There is no application server
of Yap's own. Business rules live in the database, as row-level security policies and as database
functions; Supabase provides authentication, a REST interface to the database, Realtime delivery,
file storage and scheduled jobs around it. This chapter describes both tiers and the boundary between
them. \Cref{fig:architecture} gives the overall picture.

\begin{figure}[ht]
  \centering
  % System architecture: client layers and Supabase services (Figure: architecture).
\begin{tikzpicture}[
    band/.style={box, clientbox, minimum height=8.5mm, font=\sffamily\footnotesize},
    svc/.style={box, serverbox, minimum height=11mm, text width=24mm, font=\sffamily\footnotesize},
    path/.style={font=\ttfamily\scriptsize, text=muted},
    proto/.style={font=\sffamily\scriptsize, text=ink, fill=white, inner sep=1.2pt, align=center},
  ]
  % ---------------------------------------------------------------- client
  \begin{scope}[on background layer]
    \node[layer, fit={(-0.15,0.85) (15.7,-5.35)}, inner sep=0pt] (clientframe) {};
  \end{scope}
  \node[layerlabel] at (clientframe.north west) {MOBILE CLIENT: EXPO AND REACT NATIVE ON iOS AND ANDROID};

  \node[band, minimum width=12cm] (screens) at (6,0) {Screens and navigation\enskip{\ttfamily\scriptsize src/app}};
  \node[band, minimum width=2.75cm] (components) at (1.375,-1.2) {Components\\[-2pt]{\ttfamily\scriptsize src/components}};
  \node[band, minimum width=2.75cm] (hooks) at (4.458,-1.2) {Hooks\\[-2pt]{\ttfamily\scriptsize src/hooks}};
  \node[band, minimum width=2.75cm] (theme) at (7.542,-1.2) {Theme\\[-2pt]{\ttfamily\scriptsize src/theme}};
  \node[band, minimum width=2.75cm] (lib) at (10.625,-1.2) {Pure helpers\\[-2pt]{\ttfamily\scriptsize src/lib}};
  \node[band, accentbox, minimum width=12cm] (store) at (6,-2.4) {Store: all shared state, actions and Realtime subscriptions\enskip{\ttfamily\scriptsize src/state/store.ts}};
  \node[band, accentbox, minimum width=12cm] (api) at (6,-3.6) {API layer: every backend call, error translation, signed URLs\enskip{\ttfamily\scriptsize src/data/api.ts}};
  \node[band, minimum width=12cm] (client) at (6,-4.8) {Supabase client (supabase-js), session kept in AsyncStorage\enskip{\ttfamily\scriptsize src/data/supabase-client.ts}};

  \node[box, clientbox, text width=26mm, align=left, font=\sffamily\scriptsize, minimum height=5.6cm, anchor=north west] (device) at (12.6,0.42)
    {\textbf{Device services}\\[3pt]
     Foreground location\\
     Camera and photo library\\
     Image re-encoding\\
     AsyncStorage\\
     Network state\\
     Keep awake\\
     Haptics\\[3pt]
     {\color{muted}Used by screens, hooks, the store and the API layer through Expo modules.}};

  \draw[flow] (screens.south -| components) -- (components);
  \draw[flow] (screens.south -| hooks) -- (hooks);
  \draw[flow] (hooks.south) -- (hooks |- store.north);
  \draw[flow] (6,0 |- screens.south) -- (6,0 |- store.north);
  \draw[flow] (store) -- (api);
  \draw[softflow] (screens.east) -- (12.3,0 |- screens.east) -- (12.3,0 |- api.east)
    node[pos=0.5, rotate=90, fill=white, inner sep=1pt, font=\sffamily\tiny, text=muted] {direct calls} -- (api.east);
  \draw[softflow, -] (5.5,0 |- hooks.south) -- (5.5,-1.82) -- (12.3,-1.82);
  \draw[flow] (api) -- (client);

  % ---------------------------------------------------------------- Supabase
  \begin{scope}[on background layer]
    \node[layer, fit={(-0.15,-6.85) (15.7,-11.45)}, inner sep=0pt] (serverframe) {};
  \end{scope}
  \node[layerlabel, anchor=north east] at (serverframe.north east) {SUPABASE PROJECT (HOSTED)};

  \node[svc] (rest) at (1.375,-7.95) {Data API\\[-2pt]{\scriptsize PostgREST}};
  \node[svc] (realtime) at (4.458,-7.95) {Realtime\\[-2pt]{\scriptsize Postgres changes}};
  \node[svc] (storage) at (7.542,-7.95) {Storage\\[-2pt]{\scriptsize private bucket \texttt{photos}}};
  \node[svc] (auth) at (10.625,-7.95) {Auth\\[-2pt]{\scriptsize email and password}};
  \node[svc, text width=26mm] (smtp) at (14.2,-7.95) {SMTP sender\\[-2pt]{\scriptsize Resend, mail.samarxkx.dev}};

  \node[box, serverbox, minimum width=12cm, text width=11.6cm, align=flush left, font=\sffamily\scriptsize] (pg) at (6,-10.25)
    {\textbf{\footnotesize Postgres with PostGIS}\\[2pt]
     \texttt{public}: tables with row-level security; \texttt{SECURITY DEFINER} functions for writes other than own profile, unblock and photos\\
     \texttt{private}: helper functions and \texttt{location\_changes}; not exposed through the Data API\\
     \texttt{auth}, \texttt{storage}: accounts and file metadata, managed by Supabase\\
     \texttt{pg\_cron}: the job \texttt{yap-expire-stale}, every minute};

  \draw[flow, <->] (client.south -| rest) -- node[proto] {HTTPS\\[-2pt]rpc, select} (rest);
  \draw[flow, <->] (client.south -| realtime) -- node[proto] {WebSocket} (realtime);
  \draw[flow, <->] (client.south -| storage) -- node[proto] {HTTPS\\[-2pt]upload, sign} (storage);
  \draw[flow, <->] (client.south -| auth) -- node[proto] {HTTPS\\[-2pt]sessions, codes} (auth);
  \draw[flow] (auth) -- (smtp);
  \draw[flow] (smtp.south) -- ++(0,-0.75) node[below, font=\sffamily\scriptsize, align=center, inner sep=1pt] {8-digit codes\\[-2pt]to the user's inbox};

  \draw[flow] (rest) -- node[proto] {SQL as the caller} (rest |- pg.north);
  \draw[flow] (realtime |- pg.north) -- node[proto] {row changes} (realtime);
  \draw[flow] (storage) -- node[proto] {read policy} (storage |- pg.north);
  \draw[flow] (auth) -- node[proto] {\texttt{auth.users}} (auth |- pg.north);
\end{tikzpicture}

  \caption{System architecture. Arrows inside the client show the direction of calls between layers;
    the dashed arrow is the direct calls from a few screens and one hook to the API layer, which bypass the
    store. The client reaches four Supabase services; each service reaches Postgres, where the rules are
    enforced.}
  \label{fig:architecture}
\end{figure}

\section{Technology stack}

\Cref{tab:stack} lists the main technologies and the role each plays. Versions are those pinned in
\code{package.json}.

\begin{table}[ht]
  \centering
  \caption{Technology stack.}
  \label{tab:stack}
  \tablefont
  \begin{tabularx}{\textwidth}{@{}L{34mm}L{48mm}Y@{}}
    \toprule
    \textsf{\textbf{Concern}} & \textsf{\textbf{Technology}} & \textsf{\textbf{Role in Yap}} \\
    \midrule
    App runtime & Expo SDK 57, React Native 0.86.3, React 19.2.3 & One code base for iOS and Android \\
    Language & TypeScript 6.0 in strict mode & All client code; path alias \code{@/} for \code{src/} \\
    Navigation & Expo Router 57 with typed routes & File-based routes, guarded stacks, tabs, sheets and modals \\
    Compilation & React Compiler (Expo experiment) & Automatic memoisation of components \\
    Maps & \code{react-native-maps} 1.27.2 & The \ui{Map} tab, event pins and the radius circle \\
    Animation & Reanimated 4.5, Worklets 0.10 & Press feedback, the pulsing \ui{YAP} button, the mascot \\
    Images & \code{expo-image}, \code{expo-image-picker}, \code{expo-image-manipulator} & Cached display, photo picking and camera, re-encoding before upload \\
    Device & \code{expo-location}, \code{expo-haptics}, \code{expo-keep-awake}, \code{expo-network}, \code{expo-clipboard} & Foreground location, haptics, screen-on while live, connectivity, pasting the emailed code \\
    Local storage & AsyncStorage 2.2 & The Supabase session and per-account preferences \\
    Backend client & \code{@supabase/supabase-js} 2.117 & Auth, database calls, Realtime and Storage \\
    Database & Postgres with PostGIS and \code{pg_cron} & Tables, row-level security, functions, spatial queries, housekeeping \\
    Backend services & Supabase Auth, Realtime, Storage & Sign-in and email codes, live delivery, private photos \\
    Email delivery & Custom SMTP through Resend & Sends the 8-digit confirmation and reset codes \\
    Tooling & ESLint 9 with \code{eslint-config-expo}; ImageMagick 7 & Linting; generating icons and splash images \\
    \bottomrule
  \end{tabularx}
\end{table}

\section{Client architecture}
\label{sec:client-architecture}

The client is organised in layers, each in its own directory under \code{src/}
(\cref{tab:client-layers}). Calls flow downward: screens use components, hooks and store actions; the
store calls the API layer; the API layer is the only code that uses the Supabase client.

\begin{table}[ht]
  \centering
  \caption{Client layers.}
  \label{tab:client-layers}
  \tablefont
  \begin{tabularx}{\textwidth}{@{}L{36mm}Y@{}}
    \toprule
    \textsf{\textbf{Directory}} & \textsf{\textbf{Responsibility}} \\
    \midrule
    \code{src/app} & Routes and screens. A file is a route; a \code{_layout.tsx} file is a navigator. Screens hold form state locally and read shared state through \code{useStore}. \\
    \code{src/components} & Shared components grouped by purpose: \code{ui} (buttons, cards, fields, lists, skeletons), \code{profile}, \code{live}, \code{chat}, \code{safety}, \code{events}, \code{account} and \code{mascot}. \\
    \code{src/hooks} & The shared one-second clock and countdowns, foreground location, nearby polling and nearby prompts, and lazy signing of a person's extra photos. \\
    \code{src/state} & The store (\code{store.ts}), the profile-setup draft (\code{profile-draft.tsx}) and per-account preferences saved on the phone (\code{saved-prefs.ts}). \\
    \code{src/data} & The API layer (\code{api.ts}), the single Supabase client (\code{supabase-client.ts}), the data shapes (\code{types.ts}), fixed lists and limits (\code{constants.ts}) and the types generated from the database schema (\code{database.types.ts}). \\
    \code{src/lib} & Pure helpers with no React: wording for names, genders and times, event categories, password checks, error text, image cache keys, tidying typed text, checking picked photos, and the server-corrected clock. \\
    \code{src/theme} & Design tokens (colours for light and dark mode, spacing, radii, sizes, type, elevation, motion) and the \code{useTheme} hook. \\
    \code{src/config} & Feature flags (\code{features.ts}) and the Supabase project address and publishable key (\code{supabase.ts}). \\
    \bottomrule
  \end{tabularx}
\end{table}

\paragraph{Layering rule} No screen or component imports the Supabase client. Most screens change
state only through store actions. A few screens that need no shared state call the API layer directly:
the code entry screen (to resend a code), \ui{Sign in} (to resend a confirmation code), \ui{Forgot
password} and \ui{Change password}. The hook \code{usePersonPhotos} also calls the API layer, to sign a
person's extra photos when their profile opens.

\section{Navigation and route guards}
\label{sec:navigation}

Expo Router turns the files under \code{src/app} into a navigation tree (\cref{fig:routes}). The root
layout, \code{src/app/_layout.tsx}, renders one stack with four \term{guarded} children. A guard is a
Boolean computed from the store; Expo Router allows a protected stack only while its guard is true, so
the app always lands on the correct step without explicit redirects (\cref{tab:guards}).

\begin{figure}[ht]
  \centering
  % Route tree and the guards of the root stack (Figure: routes).
\begin{tikzpicture}[
    stackbox/.style={box, accentbox, font=\sffamily\footnotesize, align=center, minimum height=11mm},
    routes/.style={box, font=\ttfamily\scriptsize, align=flush left, anchor=north, inner sep=4pt},
    guard/.style={font=\sffamily\scriptsize, text=accentdark, align=center},
  ]
  \node[box, clientbox, font=\sffamily\footnotesize] (root) at (7.9,0)
    {Root stack\enskip{\ttfamily\scriptsize src/app/\_layout.tsx}\\[-1pt]{\scriptsize exactly one guarded stack is allowed at a time}};

  \node[stackbox, text width=30mm] (auth) at (1.55,-1.75) {\texttt{(auth)}\\[-1pt]{\scriptsize signed out}};
  \node[stackbox, text width=30mm] (setup) at (5.05,-1.75) {\texttt{profile-setup}\\[-1pt]{\scriptsize signed in, no profile}};
  \node[stackbox, text width=25mm] (verify) at (8.3,-1.75) {\texttt{verify-face}\\[-1pt]{\scriptsize profile, not verified}};
  \node[stackbox, text width=50mm] (app) at (13.05,-1.75) {\texttt{(app)}\\[-1pt]{\scriptsize signed in, profile, verified or check off}};

  \foreach \n in {auth, setup, verify, app} {\draw[flow] (root.south) -- (\n.north);}

  \node[routes, text width=30mm] at (1.55,-2.65)
    {welcome\\sign-in\\sign-up\\verify {\sffamily(code)}\\forgot-password\\new-password};
  \node[routes, text width=30mm] at (5.05,-2.65)
    {index {\sffamily(18+ check)}\\name\\gender\\about\\interests\\photo\\photos};
  \node[routes, text width=25mm] at (8.3,-2.65)
    {verify-face\\[2pt]{\sffamily\color{muted} reached only when\\ the selfie check is\\ switched on}};
  \node[routes, text width=50mm] at (13.05,-2.65)
    {\begin{tabular}{@{}>{\sffamily\bfseries}l@{\enskip}>{\raggedright\arraybackslash}p{38mm}@{}}
       Tabs & index {\sffamily(Go live)}, nearby, map, profile \\[1pt]
       Sheets & go-available, person/[id], event/[id] \\[1pt]
       Modals & waiting, report/[id], incoming {\sffamily(full screen)} \\[1pt]
       Pushed & chat, event/new, group/[id], tutorial, edit-profile, settings, blocked, password \\
     \end{tabular}};
\end{tikzpicture}

  \caption{The route tree. The four children of the root stack are mutually exclusive; the guard under
    each name states when it is allowed. Not drawn: the two unguarded reading pages, \code{guidelines} and
    \code{acknowledgements}.}
  \label{fig:routes}
\end{figure}

\begin{table}[ht]
  \centering
  \caption{Root guards. \emph{Verified} is true when the profile is verified or the selfie check is
    switched off, which is the current setting.}
  \label{tab:guards}
  \tablefont
  \begin{tabularx}{\textwidth}{@{}L{30mm}L{58mm}Y@{}}
    \toprule
    \textsf{\textbf{Stack}} & \textsf{\textbf{Guard}} & \textsf{\textbf{Purpose}} \\
    \midrule
    \code{(auth)} & not signed in & Welcome, sign-up, sign-in, emailed codes, password reset \\
    \code{profile-setup} & signed in and no profile row & The 18+ confirmation and six setup steps \\
    \code{account-restricted} & signed in, profile, hidden or suspended & Explains the hold, with the support email, a sign-out and account deletion \\
    \code{verify-face} & signed in, profile, active, not verified & Selfie check (switched off) \\
    \code{(app)} & signed in, profile, active, verified & The app itself \\
    \bottomrule
  \end{tabularx}
\end{table}

Beside the guarded stacks, the root stack has two pages with no guard, \code{guidelines} (the community
guidelines) and \code{acknowledgements} (the open-source licences), so the sign-up screen and Settings can
both open them, signed in or out.

Before the store has restored the session (\code{booted} is false), the root layout renders an empty
view in the background colour. If the user is signed in but the profile could not be loaded
(\code{loadFailed}), it renders the \ui{Can't reach Yap} screen instead of any stack, so an existing user
is never sent into profile setup by a network failure (\cref{sec:offline}). An offline banner is drawn
over every screen while the phone has no connection.

The \code{(app)} layout hosts the four tabs (\ui{Go live}, \ui{Nearby}, \ui{Map}, \ui{Profile}) and every
screen opened from them. It also runs the behaviour that must apply whichever screen is open:

\begin{itemize}
  \item It starts nearby polling and nearby prompts while the user is live (\cref{sec:feature-nearby}).
  \item It sets a timer that takes the user offline when the live session reaches its end time.
  \item It pushes the \ui{incoming} screen when a new request arrives, and the \ui{chat} screen when a
    chat opens, each only once per request or chat.
  \item It opens the tutorial once after a new account saves its first profile.
  \item It shows the nearby prompt (a toast) and keeps the screen awake while the user is live.
\end{itemize}

\Cref{tab:presentations} lists how each screen in \code{(app)} is presented.

\begin{table}[ht]
  \centering
  \caption{Presentation of the screens in the \code{(app)} stack.}
  \label{tab:presentations}
  \tablefont
  \begin{tabularx}{\textwidth}{@{}L{40mm}L{40mm}Y@{}}
    \toprule
    \textsf{\textbf{Route}} & \textsf{\textbf{Presentation}} & \textsf{\textbf{Notes}} \\
    \midrule
    \code{go-available}, \code{person/[id]} & Form sheet, 85\,\% height & Grabber visible \\
    \code{event/[id]} & Form sheet, 75\,\% height & Event details, \ui{Join}; \ui{Delete event} for the host \\
    \code{waiting}, \code{report/[id]} & Modal, slides from the bottom & \\
    \code{incoming} & Full-screen modal & Swipe to dismiss disabled; the user must answer or let it expire \\
    \code{tutorial} & Pushed, fades in & Swipe back disabled \\
    \code{chat}, \code{event/new}, \code{group/[id]}, \code{edit-profile}, \code{settings}, \code{blocked}, \code{password} & Pushed screen & Slides from the right on both platforms \\
    \bottomrule
  \end{tabularx}
\end{table}

\section{State management}
\label{sec:state}

\subsection{The store}

All shared state lives in one module, \code{src/state/store.ts} (\cref{fig:store-flow}). The module holds
a single immutable state object and a set of listeners. A private function, \code{set(patch)}, replaces
the object with a merged copy and notifies every listener. Components read state with
\code{useStore(selector)}, built on React's \code{useSyncExternalStore}, so a component re-renders only
when the slice it selects changes. Nothing outside the module writes state: screens call exported actions,
and the store's own listeners, timers and Realtime handlers call the same actions.
\Cref{app:store} lists every state field and action.

\begin{figure}[ht]
  \centering
  % Store data flow: what triggers actions, what they change, and who reads it (Figure: store flow).
\begin{tikzpicture}[
    inbox/.style={box, clientbox, font=\sffamily\scriptsize, text width=24mm, minimum height=9mm, align=center},
    col/.style={box, font=\sffamily\scriptsize, align=flush left, anchor=north, inner sep=5pt, minimum height=56mm},
    lab/.style={font=\sffamily\scriptsize, text=muted, fill=white, inner sep=1.2pt, align=center},
  ]
  \node[layerlabel, inner sep=0pt, anchor=west] at (0.05,0.75) {TRIGGERS};
  \node[inbox] (ui) at (1.4,0) {Screens: user actions};
  \node[inbox] (timers) at (1.4,-1.2) {Timers: session end,\\answer timeout};
  \node[inbox] (listeners) at (1.4,-2.4) {App-state and\\network listeners};
  \node[inbox] (authl) at (1.4,-3.6) {Auth listener};
  \node[inbox, accentbox] (rt) at (1.4,-4.8) {Realtime handlers};

  \node[col, accentbox, text width=36mm] (actions) at (5.35,0.45)
    {\textbf{\footnotesize Store actions}\\[4pt]
     \textbf{session}\enskip\texttt{init}, \texttt{signOut}\\[3pt]
     \textbf{live}\enskip\texttt{goGreen}, \texttt{goOffline}\\[3pt]
     \textbf{requests}\enskip\texttt{sayHi}, \texttt{cancelRequest}\\[3pt]
     \textbf{chat}\enskip\texttt{sendMessage}, \texttt{endChat}\\[3pt]
     \textbf{safety}\enskip\texttt{blockUser}\\[3pt]
     \textbf{events}\enskip\texttt{loadEvents}, \texttt{joinEvent}\\[3pt]
     \textbf{recovery}\enskip\texttt{catchUp}\\[5pt]
     {\color{muted}full list: Appendix B}};

  \node[col, text width=32mm] (state) at (9.9,0.45)
    {\textbf{\footnotesize State: one object}\\{\color{muted}replaced by \texttt{set(patch)}}\\[4pt]
     \textbf{account}\enskip\texttt{session}, \texttt{profile}\\[3pt]
     \textbf{live}\enskip\texttt{availability}, \texttt{nearby}\\[3pt]
     \textbf{requests}\enskip\texttt{outgoing}, \texttt{incoming}\\[3pt]
     \textbf{chat}\enskip\texttt{chat}\\[3pt]
     \textbf{safety}\enskip\texttt{blockedIds}\\[3pt]
     \textbf{events}\enskip\texttt{events}\\[3pt]
     \textbf{flags}\enskip\texttt{booted}, \texttt{loadFailed}};

  \node[col, clientbox, text width=20mm, align=center] (screens) at (14.3,0.45)
    {\textbf{\footnotesize Screens}\\[4pt]read slices with \texttt{useStore}\\[4pt]re-render only when their slice changes};

  \node[box, accentbox, text width=36mm, font=\sffamily\scriptsize] (api) at (5.35,-7.0) {API layer \texttt{api.ts}};
  \node[box, serverbox, text width=32mm, font=\sffamily\scriptsize] (sb) at (9.9,-7.0) {Supabase: rpc, select,\\storage, auth, Realtime};

  \foreach \n in {ui, timers, listeners, authl, rt} {\draw[flow] (\n.east) -- (\n.east -| actions.west);}
  \draw[flow] (actions.east |- 0,-1.6) -- (state.west |- 0,-1.6);
  \draw[flow] (state.east |- 0,-1.6) -- node[lab, above=1pt] {notify} (screens.west |- 0,-1.6);
  \draw[flow, <->] (actions.south) -- node[lab, right=3pt, fill=none, align=left] {calls; results, or refusals\\as sentences} (api.north);
  \draw[flow, <->] (api) -- (sb);
  \draw[push] (sb.south) -- ++(0,-0.55) -| node[lab, pos=0.25, below=1pt] {pushed inserts and updates on requests and messages} (rt.south);
\end{tikzpicture}

  \caption{Data flow through the store. Five kinds of trigger call actions; actions call the API layer and
    update the one state object; screens re-render from the slices they select. Realtime changes arrive
    as handler calls, like any other trigger.}
  \label{fig:store-flow}
\end{figure}

\subsection{State outside the render tree}

Some state never needs to re-render anything and is kept in module variables inside the store. It
records work in progress and guards against races (\cref{sec:ordering}):

\begin{itemize}
  \item the unsubscribe functions of the three Realtime channels the store may hold open: incoming
    requests, the one-to-one chat, and the open group chat;
  \item \code{handledRequestIds} and \code{closedChatIds}, which stop a catch-up from bringing back a
    request already answered or a chat already closed on this phone;
  \item three counters, \code{nearbyRefreshes}, \code{eventLoads} and \code{eventChatVisits}, which let
    the store discard an answer that arrives after a newer request has started;
  \item flags that coalesce or order work: \code{catchUpRunning}, \code{catchUpAgain},
    \code{offlineRetry}, \code{watchingId}, \code{openingChatId}, \code{extrasLoaded} and
    \code{initStarted}.
\end{itemize}

\subsection{Other client state}

Profile setup keeps its answers in a React context, \code{ProfileDraftProvider}, so the user can move
back and forth between the six steps without losing input. Two preferences are saved on the phone per
account under the AsyncStorage key \code{yap.prefs.<user id>}: how many times the chat safety tip has
been shown, and the map's event radius. Forms keep their field values in component state until they
are submitted.

\section{The API layer}
\label{sec:api-layer}

\code{src/data/api.ts} is the client's only route to the backend. It has four responsibilities.

\begin{description}
  \item[Wrapping calls.] Each exported function wraps one Supabase call or a short sequence of them:
    an Auth method, a database function through \code{supabase.rpc}, a read of a table the user may
    see, or a Storage operation. It converts the JSON the server returns into the client's types in
    \code{types.ts}.
  \item[Translating errors.] A database function refuses a request by raising an exception with a short
    fixed message such as \code{not nearby}. The API layer maps every such message to a sentence for
    the screen through the \code{SERVER_REFUSALS} table, maps Auth error codes in the same way, and
    replaces anything else with a generic sentence, so raw database text never reaches the user
    (\cref{app:refusals}).
  \item[Realtime.] It builds the Realtime channels and subscribes them as the signed-in user
    (\cref{sec:realtime}).
  \item[Photos.] It re-encodes and uploads photos, removes unused files, and turns storage paths into
    signed URLs through a cache (\cref{sec:flow-photos}).
\end{description}

\section{Backend architecture}
\label{sec:backend-architecture}

The backend is one hosted Supabase project. The client knows its URL and its publishable key, both of
which are public by design: the key grants only what row-level security and the database functions
allow. No secret key ships in the app.

\paragraph{Schemas} Yap's database objects span six schemas. \code{public} holds the app's tables and the
database functions the client calls; it is the schema the Data API exposes. \code{private} holds helper
functions and the table \code{location_changes}; it is not exposed. \code{auth} and \code{storage} are
managed by Supabase and hold accounts and file metadata. \code{extensions} holds PostGIS, and
\code{cron} holds the \code{pg_cron} schedule.

\paragraph{Access model} Row-level security is enabled on every application table and denies by default:
a table without a policy cannot be read or written by a client at all. Direct access is granted only
where the client must read rows itself (its own profile, the requests and messages it takes part in, its
own block list) and where a write is trivially safe (its own profile, removing its own block). Every
other write is a \code{SECURITY DEFINER} function: it runs with the privileges of its owner, reads the
caller's identity from \code{auth.uid()}, and checks the rules before changing anything. \Cref{ch:server-interface}
lists the functions and \cref{ch:security} the policies.

\paragraph{Realtime} The tables \code{chat_requests} and \code{messages} are in the
\code{supabase_realtime} publication with full replica identity. Realtime delivers a change to a
subscriber only if the subscriber's token may read the row under the table's select policy, so each user
receives only the requests and messages they take part in.

\paragraph{Storage} Photos live in a private bucket, \code{photos}, under a folder named after the
owner's user ID. Clients write only inside their own folder. Reading, and with it signing a URL, is
allowed only for photos of someone the reader can currently see (\cref{sec:photo-privacy}).

\paragraph{Scheduled work} A \code{pg_cron} job named \code{yap-expire-stale} runs
\code{private.expire_stale()} every minute. It ends live sessions whose time is up, expires unanswered
requests, ends chats in which neither person is live, and prunes the location-change log
(\cref{sec:housekeeping}).

\paragraph{Authentication and email} Supabase Auth provides email and password accounts. \ui{Confirm
email} is switched on, the minimum password length is 8, and the email templates for sign-up
confirmation and password reset show an 8-digit code instead of a link, because the app does not open
email links. Auth sends these emails through a custom SMTP sender, Resend, on the domain
\code{mail.samarxkx.dev}. These settings are made in the Supabase dashboard (\cref{sec:dashboard}).
